\documentclass[11pt]{article}

\usepackage[margin=1in]{geometry}
\usepackage{amsmath,amssymb,amsthm}
\usepackage{graphicx}
\usepackage[hidelinks]{hyperref}

\newtheorem{theorem}{Theorem}

\newcommand{\X}{\mathfrak{X}_{0,1}^{n}}
\newcommand{\SSing}{\mathcal{S}}

\title{An Adjoint-Strictly-Singular Subcase of the Dual Problem}
\author{Run \texttt{fa\_banach\_001}, agent \texttt{agent\_lane\_09}}
\date{June 15, 2026}

\begin{document}
\maketitle

\section*{Status}

This packet records a \texttt{partial\_result\_likely\_valid} for
Problem 2(iii) in Argyros--Beanland--Motakis
\cite{ABM}.  The result is not a full solution of the dual problem; it
handles the subcase in which the dual operators have strictly singular
adjoints on the reflexive predual.

\section*{Source Problem}

The source paper constructs, for each $n\in\mathbb N$, a reflexive
Tsirelson-like space $\X$ such that for every infinite-dimensional subspace
$Y$ of $\X$, every product of $n+1$ strictly singular operators on $Y$ is
compact.  In its final ``Problems and Questions'' section, the paper asks:

\begin{quote}
Does $\mathfrak X_{0,1}^{n*}$ satisfy that whenever
$S_1,\ldots,S_{n+1}:\mathfrak X_{0,1}^{n*}\to
\mathfrak X_{0,1}^{n*}$ are strictly singular, then the composition
$S_1\cdots S_{n+1}$ is compact, as in $\mathfrak X_{0,1}^n$?
\end{quote}

The question appears on PDF page 43 of \cite{ABM}; the local source lines are
approximately \texttt{source.tex:3867--3874}.

\begin{center}
\includegraphics[width=\textwidth]{figures/open_problem_crop.png}
\end{center}

\section*{Proof Intuition}

The source theorem already proves the desired product-compactness property on
the predual $\X$.  Since $\X$ is reflexive, every operator on ${\X}^*$ has an
adjoint that can be identified with an operator on $\X$.  Therefore the dual
problem would follow immediately for any family whose adjoints are strictly
singular on $\X$: apply the source theorem to the adjoints, then use
Schauder's theorem to pass compactness back to the original product.

This isolates the missing point in the full problem.  Strict singularity is
not self-dual for arbitrary Banach spaces, so one cannot simply assume that a
strictly singular operator on ${\X}^*$ has strictly singular adjoint on $\X$.

\section*{The Partial Result}

\begin{theorem}
Let $n\in\mathbb N$ and let $\X=\mathfrak X_{0,1}^n$ be the reflexive space
constructed in \cite{ABM}.  Let
$S_1,\ldots,S_{n+1}\in\mathcal L({\X}^*)$.  Assume that, after identifying
${\X}^{**}$ with $\X$, each adjoint
\[
   S_i^*:\X\longrightarrow \X
\]
is strictly singular.  Then $S_1S_2\cdots S_{n+1}$ is compact on ${\X}^*$.
\end{theorem}

\begin{proof}
By the main theorem of \cite{ABM}, applied to the subspace $Y=\X$, the
composition of any $n+1$ strictly singular operators on $\X$ is compact.
Since each $S_i^*$ is strictly singular by hypothesis, the operator
\[
   S_{n+1}^*S_n^*\cdots S_1^*:\X\longrightarrow \X
\]
is compact.

The adjoint of the original product is
\[
   (S_1S_2\cdots S_{n+1})^*
      = S_{n+1}^*S_n^*\cdots S_1^*,
\]
and hence is compact.  By Schauder's theorem, a bounded operator is compact
if and only if its adjoint is compact.  Therefore
$S_1S_2\cdots S_{n+1}$ is compact on ${\X}^*$.
\end{proof}

\section*{Equivalent Reduction}

Because $\X$ is reflexive, every bounded operator on ${\X}^*$ has a preadjoint
on $\X$: namely, for $S\in\mathcal L({\X}^*)$, identify $S^*$ as an operator
from ${\X}^{**}$ to ${\X}^{**}$, hence from $\X$ to $\X$.  Thus Problem 2(iii)
would be solved if one could prove the implication
\[
   S\in\SSing({\X}^*) \quad\Longrightarrow\quad S^*\in\SSing(\X).
\]
The theorem above proves the dual product conclusion exactly under this
additional adjoint-strict-singularity hypothesis.

\section*{Limitations and Novelty Check}

This packet does not prove the implication
$S\in\SSing({\X}^*)\Rightarrow S^*\in\SSing(\X)$.  That implication is false for
general Banach spaces, so it would require a special property of
$\mathfrak X_{0,1}^n$ or its dual.

Local run indexes were searched for arXiv:1309.4516, the paper title,
``dual strictly singular'', and adjoint/strictly singular keywords.  No
duplicate packet was found.  Web searches on June 15, 2026 found the source
paper, the related higher-order distortion paper arXiv:1408.5065, and the
related asymptotic-symmetry paper arXiv:1902.10098, but no later full answer
to Problem 2(iii).

\section*{Human Review Recommendation}

Review as a small partial subcase/reduction.  The proof itself is short; the
main point to verify is that the statement is not overread as a full solution.
The remaining full problem is whether strict singularity on ${\X}^*$ forces
strict singularity of the adjoint on $\X$ for this particular reflexive
Tsirelson-like space.

\begin{thebibliography}{9}

\bibitem{ABM}
S. A. Argyros, K. Beanland, and P. Motakis,
\emph{Strictly singular operators in Tsirelson like spaces},
arXiv:1309.4516.

\bibitem{Schauder}
J. Schauder's compactness theorem: a bounded operator between Banach spaces is
compact if and only if its adjoint is compact.

\end{thebibliography}

\end{document}
